\section{Sparse Tensor and Convolution}
\label{sec:sparse}

In traditional speech, text, or image data, features are extracted densely. Thus, the most common representations of these data are vectors, matrices, and tensors. However, for 3-dimensional scans or even higher-dimensional spaces, such dense representations are inefficient due to the sparsity. Instead, we can only save the non-empty part of the space as its coordinates and the associated features. This representation is an N-dimensional extension of a sparse matrix; thus it is known as a sparse tensor. There are many ways to save such sparse tensors compactly~\cite{sparsetensor}, but we follow the COO format as it is efficient for neighborhood queries (Sec.~\ref{sec:sparse_convolution}). Unlike the traditional sparse tensors, we augment the sparse tensor coordinates with the batch indices to distinguish points that occupy the same coordinate in different batches~\cite{sparse3dcnn}. Concisely, we can represent a set of 4D coordinates as $\mathcal{C} = \{(x_i, y_i, z_i, t_i)\}_i$ or as a matrix $C$ and a set of associated features $\mathcal{F} = \{\mbf{f}_i\}_i$ or as a matrix $F$. Then, a sparse tensor can be written as
\begin{equation}
C = \begin{bmatrix}
x_1 & y_1 & z_1 & t_1 & b_1\\ 
    &     & \vdots &  &     \\ 
x_N & y_N & z_N & t_N & b_N
\end{bmatrix}, \;
F = \begin{bmatrix}
\mathbf{f}_1^T\\ 
\vdots\\ 
\mathbf{f}_N^T
\end{bmatrix}
\end{equation}
where $b_i$ and $\mathbf{f}_i$ are the batch index and the feature associated to the $i$-th coordinate. In Sec.~\ref{sec:crf}, we augment the 4D space with the 3D chromatic space and create a 7D sparse tensor for \textit{trilateral} filtering.


\subsection{Generalized Sparse Convolution}
\label{sec:sparse_convolution}

In this section, we generalize the sparse convolution~\cite{sparsecnn, sparse3dcnn} for generic input and output coordinates and for arbitrary kernel shapes. The generalized sparse convolution encompasses not only all sparse convolutions but also the conventional dense convolutions. Let $x^{\text{in}}_\mbf{u} \in \mathbb{R}^{N^\text{in}}$ be an $N^\text{in}$-dimensional input feature vector in a $D$-dimensional space at $\mbf{u} \in \mathbb{R}^D$ (a D-dimensional coordinate), and convolution kernel weights be $\mathbf{W} \in \mathbb{R}^{K^D
\times N^\text{out} \times N^\text{in}}$. We break down the weights into spatial weights with $K^D$ matrices of size $N^\text{out} \times N^\text{in}$ as $W_\mbf{i}$ for $|\{\mbf{i}\}_\mbf{i}| = K^D$. Then, the conventional dense convolution in D-dimension is
\begin{equation}
\mathbf{x}^{\text{out}}_\mbf{u} = \sum_{\mbf{i} \in \mathcal{V}^D(K)} W_\mbf{i} \mathbf{x}^{\text{in}}_{\mbf{u} + \mbf{i}} \text{ for } \mbf{u} \in \mathbb{Z}^D,
\label{eq:dense_convolution}
\end{equation}
where $\mathcal{V}^D(K)$ is the list of offsets in D-dimensional hypercube centered at the origin. e.g. $\mathcal{V}^1(3)=\{-1, 0, 1\}$. The generalized sparse convolution in Eq.~\ref{eq:sparse_convolution} relaxes Eq.~\ref{eq:dense_convolution}.
\begin{equation}
    \mathbf{x}^{\text{out}}_\mbf{u} = \sum_{\mbf{i} \in \mathcal{N}^D(\mbf{u}, \mathcal{C}^{\text{in}})} W_\mbf{i} \mathbf{x}^{\text{in}}_{\mbf{u} + \mbf{i}} \text{ for } \mbf{u} \in \mathcal{C}^{\text{out}}
\label{eq:sparse_convolution}
\end{equation}
where $\mathcal{N}^D$ is a set of offsets that define the shape of a kernel and $\mathcal{N}^D(\mbf{u}, \mathcal{C}^\text{in})= \{\mbf{i} | \mbf{u} + \mbf{i} \in \mathcal{C}^\text{in}, \mbf{i} \in \mathcal{N}^D \}$ as the set of offsets from the current center, $\mbf{u}$, that exist in $\mathcal{C}^\text{in}$. $\mathcal{C}^\text{in}$ and $\mathcal{C}^\text{out}$ are predefined input and output coordinates of sparse tensors. First, note that the input coordinates and output coordinates are not necessarily the same. Second, we define the shape of the convolution kernel arbitrarily with $\mathcal{N}^D$. This generalization encompasses many special cases such as the dilated convolution and typical hypercubic kernels. Another interesting special case is the "sparse submanifold convolution"~\cite{sparse3dcnn} when we set $\mathcal{C}^\text{out} = \mathcal{C}^\text{in}$ and $\mathcal{N}^D = \mathcal{V}^D(K)$. If we set $\mathcal{C}^\text{in} = \mathcal{C}^\text{out} = \mathbb{Z}^D$ and $\mathcal{N}^D = \mathcal{V}^D(K)$, the generalized sparse convolution becomes the conventional dense convolution (Eq.~\ref{eq:dense_convolution}). If we define the $\mathcal{C}^\text{in}$ and $\mathcal{C}^\text{out}$ as multiples of a natural number and $\mathcal{N}^D = \mathcal{V}^D(K)$, we have a strided dense convolution.
